% v3_draft.tex — Ceremonial Entheogen Use in Puerto Rico
% Phase 5 (peer review) output. Applies every accepted finding from
% manuscript/PEER_REVIEW.md to v2_draft.tex; see drafts/revision_notes.md for the
% item-by-item disposition. Built on manuscript/drafts/v0_jpd_template.tex.
% Numbers trace to outputs/report_v2.md except where the report does not carry a
% univariate distribution; those are direct recomputations from
% data/results_7_15_2026.xlsx via analysis/data_prep.load(), and are listed in
% revision_notes.md.
% Formatting per manuscript/JOURNAL_REQUIREMENTS.md (Journal of Psychoactive Drugs).
% Word budget: Introduction-Methods-Results-Discussion <= 4,000 words (strict, per author decision).

\documentclass[12pt]{article}

\usepackage[utf8]{inputenc}
\usepackage[T1]{fontenc}
\usepackage{times}
\usepackage[margin=1in]{geometry}
\usepackage{setspace}
\doublespacing
\usepackage{lineno}
\usepackage{graphicx}
\usepackage{booktabs}
\usepackage{longtable}  % Table 3 gained structural-skip rows in v3 and no longer fits one page
\usepackage{amsmath}
\usepackage{url}
\urlstyle{same}  % render URLs in the body font; avoids a Courier dependency
\emergencystretch=3em  % double spacing leaves little slack; prevents overfull lines at long citations
\usepackage[authoryear,round]{natbib}
\bibliographystyle{chicago}

\title{Ceremonial Entheogen Use in Puerto Rico: A Descriptive Survey of Participants
and a Small Facilitator Subsample}

\author{
  Yamil O. Ortiz Ortiz$^{1}$ \and Jean C. V\'elez Rodr\'iguez$^{2}$ \and Juli\'an M. Hern\'andez Torres$^{2,3}$ \and
  Juliana Mill\'an-Torres$^{4}$ \and Paulina Rull\'an Farinacci$^{5}$ \and Carolina Riollano$^{6}$ \and
  Ang\'elica N. Grajales Rold\'an$^{1}$ \and Adriana I. Rodr\'iguez Massa$^{1}$ \\[1em]
  \begin{minipage}{0.92\textwidth}\centering\small\setstretch{1.15}
  $^{1}$Centro de Investigaciones Sociales, Facultad de Ciencias Sociales, University of Puerto Rico,
  R\'io Piedras Campus, San Juan, Puerto Rico \\
  $^{2}$Puerto Rico Institute for Psychedelic Science, Medicine \& Awareness, San Juan, Puerto Rico \\
  $^{3}$Behavioral Sciences Research Institute, University of Puerto Rico, Medical Sciences Campus,
  San Juan, Puerto Rico \\
  $^{4}$Colectivo Psicod\'elico de Puerto Rico, Colectivo Psicod\'elico Inc., and Within Psychological
  Services, San Juan, Puerto Rico \\
  $^{5}$Department of Psychiatry, School of Medicine, University of Puerto Rico, Medical Sciences Campus,
  San Juan, Puerto Rico \\
  $^{6}$Colectivo Psicod\'elico de Puerto Rico, San Juan, Puerto Rico \\[0.5em]
  Corresponding author: Jean C. V\'elez Rodr\'iguez, jean.velez5@upr.edu
  \end{minipage}
}

\date{}

\begin{document}

\maketitle
\linenumbers

\begin{abstract}
\noindent
% ABSTRACT — unstructured, <=200 words (JPD requirement).
Ceremonial entheogen use occurs in Puerto Rico outside clinical and regulatory structures, and the
practices surrounding it have not been described. We conducted an anonymous, bilingual, cross-sectional
online survey of adults aged 21 and over reporting involvement in entheogen-assisted ceremonies in Puerto
Rico, recruited by convenience and chain referral, February--June 2026. Of 100 submissions, 73
identified a role: 72 as participants and 8 as facilitators, 7 of whom were both. Respondents were
middle-aged (mean 42.4 years), mostly university-educated, and almost all resident in Puerto Rico. Most
recent ceremonies were typically single, multi-hour events in natural settings involving psilocybin
mushrooms or ayahuasca, and satisfaction was strongly ceiling-weighted (median 6 of 6). Preparation was
near-universal among those answering while pre-ceremony screening was not, and the small facilitator
subsample described training and emergency planning as informal rather than credentialed. Five of 67
participants answering reported ever experiencing non-consensual physical contact during a ceremony, and 2
of 65 at their most recent one. A comparison of any preparation against none, specified in the analysis
plan before testing, was positive in direction but rests on five unprepared respondents and is exploratory.
These findings describe this sample only and support no prevalence or causal claim.
\end{abstract}

\noindent\textbf{Keywords:} entheogens; psychedelics; ceremonial use; Puerto Rico; harm reduction; survey research

\section{Introduction}

Ceremonial use of entheogens---psilocybin-containing mushrooms, ayahuasca, and related preparations taken
in structured group or guided settings---has become substantially more visible across the Americas over the
past two decades, largely outside the clinical-trial infrastructure that generates most of the contemporary
research literature \citep{golden2022effects,carvalho2025scoping}. The resulting evidence base is
asymmetric. Randomized clinical research characterizes carefully screened participants receiving
standardized doses under monitored conditions, and its samples remain narrow
\citep{hughes2024ethnoracial,marcus2026psychedelics}, while community ceremonial practice is documented
mainly through ethnographic and qualitative work, with comparatively few survey-level descriptions of who
takes part and under what conditions \citep{carvalho2025scoping}.

Puerto Rico occupies an unusual position in this landscape. It is a Spanish-speaking Caribbean jurisdiction
in which psilocybin, DMT, and related compounds are Schedule I substances under the federal Controlled
Substances Act \citep{uscongress1970csa} and the territory's own aligned statute \citep{pr1971csa}, and
which sits outside the state-level decriminalization and regulated-access reforms reshaping the legal
environment in parts of the United States \citep{siegel2023psychedelic}. It is also a setting with a
long-documented tradition of spiritually framed healing practice---espiritismo and related
practices---described in the clinical and public health literature for more than four decades
\citep{koss1980therapist,comasdiaz1981puertorican,bird1981sociopsychiatry,zerrate2022espiritismo}.

Local survey evidence on entheogen use is nonetheless minimal. One published Puerto Rico--specific survey
exists, conducted by members of the present research team, but it is narrower in scope: it characterizes
psilocybin use patterns, motivations, and personality correlates rather than ceremonial practice, and
addresses neither facilitation nor safety \citep{velezrodriguez2026psilocybin}. An unpublished thesis has
examined university-community perceptions of psilocybin mushrooms \citep{doeltermorales2020hongos}. Neither
describes how ceremonies are conducted. Basic descriptive questions---who takes part, which substances are
used, where and for how long ceremonies occur, whether participants are prepared or screened beforehand,
and what adverse experiences are reported---therefore have no published local answer.

A descriptive characterization contributes in three ways. First, it establishes whether practices reported
in a community sample resemble the harm-reduction structures clinical and retreat-based frameworks assume:
preparation, pre-ceremony screening, and a plan for difficult experiences
\citep{gorman2021psychedelic,pilecki2021ethical,dutton2025harm}. Second, it identifies which safety
outcomes---including boundary violations, shown elsewhere to be measurable in naturalistic samples
\citep{kruger2025psychedelic,peluso2020reflections}---occur at frequencies justifying measurement in a
larger study. Third, it produces the instrument-level evidence needed before any confirmatory study is
designed here. None of this depends on a test reaching significance: the description and its measurement
lessons are the result.

Our objectives were to describe the sample and the ceremonial context, and to examine bivariate
associations between preparation and reported ceremony outcomes and between participant-observable
facilitator practices and perceived safety. A safety and adverse-event domain, absent from the original
protocol, was added during analysis. Every objective is descriptive. The design supports no statement about
how common these practices are in Puerto Rico, and no statement about what any practice produces.

\section{Methods}

\subsection{Design, ethics, and consent}

This was a single-timepoint, cross-sectional, self-administered online survey of adults reporting
involvement in entheogen-assisted spiritual ceremonies in Puerto Rico. The sample was non-probability and
the analysis unweighted. The study was approved by the Institutional Review Board of the University of
Puerto Rico, Medical Sciences Campus (protocol 2509472552, approved 4 November 2025) and was fielded on
KoboToolbox between 24 February and 9 June 2026. Participation was fully anonymous: no names, dates of
birth, or other personal identifiers were collected. The information sheet---covering purpose, procedures,
risks, voluntariness, and data retention---was presented within the questionnaire itself, and respondents
who read it and proceeded to the items thereby gave informed consent. This procedure was described in the
protocol submitted for ethical review and approved as described, so no signed consent document was
collected or required. Every item offered ``prefer not to answer'', withdrawal was permitted at any point,
no compensation was offered, and mental-health resources were listed. Reporting follows STROBE (checklist
supplied as Supplementary Material).

\subsection{Recruitment and sample construction}

Recruitment used convenience sampling with chain referral, promoted through Facebook, Instagram, Twitter/X,
and WhatsApp and by bulletin-board announcements at the University of Puerto Rico R\'io Piedras and Medical
Sciences campuses. Inclusion criteria were age 21 or older, ability to read and write in Spanish, access to
an internet-connected device, and prior participation in an entheogen-assisted spiritual ceremony in Puerto
Rico. These criteria were self-assessed and advisory rather than enforced by instrument logic: only age
gated the questionnaire, the form was fielded bilingually, respondents reporting no ceremony involvement
reached the role fork regardless, and residence in Puerto Rico was never required. The recruitment target
of 100 responses was met.

Of the 100 submissions collected, 13 were abandoned before any substantive item and excluded, leaving 87
respondents who reached the role-identification questions. Of those, 14 answered ``no'' to both role items,
so 16\% of those reaching the role fork screened out; they are not described as a subgroup. The analytic
sample is the remaining $N=73$ respondents who identified a role. Roles are not a partition: 72 identified
as ceremony participants and 8 as facilitators, 7 of the 8 facilitators also identifying as participants,
so the facilitator branch is largely a subset of the participant branch rather than a separate group of
people. The branches are nonetheless unlinked \emph{at the ceremony level}---the instrument carries no
ceremony or facilitator identifier, and for a dual-role respondent the ceremonies led are different events
from those attended---so facilitator- and participant-reported measures cannot be crossed. One further
routing rule governs every participant-side base: dual-role respondents reached the participant block only
after a continuation prompt, and 3 of the 7 answered ``no'' and were shown no participant item. The
participant branch is therefore 72 self-identified respondents of whom 69 were routed, those 3 appearing as
structural skips throughout (Figure~\ref{fig:flow}).

\subsection{Instrument}

No validated Puerto Rico--specific instrument for characterizing ceremonial entheogen use exists, so the
questionnaire was developed by the research team, reviewed iteratively, professionally translated, and
given a final review for cross-language consistency. It was deployed as one bilingual XLSForm with Spanish
as the default and English available through the platform's language switcher, so open-text responses were
mixed-language.

The primary outcome, satisfaction with the most recent ceremony, asked ``From 0 to 6, how pleasant was the
ceremony?'' with anchors from \emph{very unpleasant} (0) to \emph{very pleasant} (6), and was treated as
ordinal. The primary exposure asked whether, aside from the facilitator's instructions, the respondent
personally engages in preparation for these ceremonies (yes / sometimes / no), capturing self-directed
preparation only. The two differ in referent: the outcome asks about one event and the exposure about a
general practice, so the primary analysis relates a single-ceremony outcome to a habitual exposure.
Secondary participant-side measures---in-ceremony comfort and safety, trust in the facilitator, prior
knowledge of the facilitator's background, whether screening questions were asked, and prior disclosure to
a health professional---all refer to the most recent ceremony. The non-consensual contact item asked
whether the respondent had \emph{ever} experienced physical touch during a ceremony that felt unexpected or
that they had not clearly agreed to beforehand, explicitly including moments of uncertainty or absent
recall; its breadth exceeds a narrower assault framing and the resulting figure should be read accordingly.
A parallel item asked the same of the most recent ceremony alone, and both are reported. That lifetime
referent does not match the facilitator-practice measures crossed against it in \S3.4, a mismatch restated
where those results appear. Facilitator-branch items covered training, screening approach,
distress-response protocol, emergency planning, and crisis signs witnessed. Three substance variables were
kept distinct and never merged: facilitator-led substances, participant substances across all ceremonies
attended, and the substance used in the most recent ceremony. Attention checks were branch-specific.
Verbatim wording and anchors appear in Supplementary Table~S1.

\subsection{Missing data, structural skips, and disclosure}

Denominators shift across items for two distinct reasons, kept separate throughout: structural skips arise
from branch logic and are reported as their own row in Tables~\ref{tab:context} and \ref{tab:safety}, while
item nonresponse is the residual. Every figure carries its own base. ``Prefer not to answer'' and ``I'm not
sure'' were retained as substantive categories in all frequency tables, excluded from ordinal numeric tests
because they carry no rank position, and treated as item-missing rather than scored as zero in composite
indices.

Counts are reported exactly throughout, including cells below five. The disclosure protection this study
relies on is its design rather than cell masking: no direct identifiers were collected, all reporting is
aggregate, and no free-text response is quoted anywhere. A masking convention would have been cosmetic in
any case, because every table states the base against which its percentages are computed, from which any
masked count is recoverable.

Because the platform populated no duplicate-detection metadata and recorded no IP or device information,
duplicate submissions were screened manually on the respondent-chosen nickname and password; no analytic
row required exclusion on duplicate grounds, and the two respondents who failed an attention check were
flagged rather than dropped, per protocol (Supplementary Material \S S7).

\subsection{Statistical approach}

The outcome distribution was strongly ceiling-weighted and several exposures had thin or unbalanced
categories, so all tests were nonparametric: Mann--Whitney $U$, Kruskal--Wallis $H$, permutation $\chi^2$
with 10,000 label permutations (seed 20260715) for thin contingency tables, and Fisher's exact test at very
small $n$. Effect sizes were rank-biserial $r_{rb}$ with the common-language effect size, epsilon-squared
$\varepsilon^2$, and Cram\'er's $V$ respectively, with Wilson 95\% confidence intervals for prevalences.

The preparation-to-satisfaction comparison was designated in the analysis plan as the primary comparison
before any test was run, and is reported uncorrected, with a tie-aware permutation test as a robustness
check because the outcome is heavily tied. The study was not preregistered, and the analysis plan was
written after data collection. All other bivariate tests were assigned to four prespecified
families---bivariate, substance-stratified, safety-domain, and legal-knowledge---corrected jointly within
family by the Benjamini--Hochberg procedure and kept separate because they share neither exposures nor
outcome domains (Supplementary Tables~S4--S7). One exploratory proportional-odds model was fitted
(Supplementary Table~S8). Analyses were conducted in Python; the full pipeline is available (see Data
availability).

\subsection{Protocol deviations}

Three deviations are reported transparently. The protocol's four-arm preparation $\times$ integration
design is not estimable from the collected data: the instrument contains no clean participant
integration-received item, and the preparation split is badly unbalanced. Preparation was therefore
examined against the outcome directly, and no integration measure is reconstructed anywhere in this paper.
The protocol also specified crossing facilitator training and protocol items against participant-perceived
safety, and a composite harm-reduction score mixing the two branches; because facilitator self-report
cannot be linked to the ceremony any particular participant attended, both would have required pairing
unrelated events. Two participant-answerable proxies were substituted---prior knowledge of the
facilitator's background and whether screening questions were asked---facilitator self-report is presented
separately and descriptively, and two branch-specific composites replaced the mixed one. Two skip-logic
errors in the instrument were also identified (Supplementary Material \S S6).

\section{Results}

\subsection{Respondents}

The analytic sample of 73 adults skewed middle-aged and highly educated (Table~\ref{tab:demo}). Among the
72 with a valid age, the mean was 42.4 years and the median 41 (SD 11.9, range 21--71, IQR 34--50); 35 of
the 71 answering the gender item (49\%) identified as women; 56 of the 72 answering the education item
(78\%) held a bachelor's degree or above; and 44 (61\%) described their economic situation as comfortable
against 6 (8\%) as difficult. Sixty-eight (94\% of the 72 answering) reported residing in Puerto Rico;
residency was a descriptive covariate rather than an inclusion filter, and the presence of respondents
living elsewhere tempers any reading of the sample as describing Puerto Rico specifically. Free-text
engagement was substantial relative to sample size: 37 open-text fields yielded 285 non-empty answers,
roughly three-quarters in Spanish. A coding scaffold for the research team's thematic analysis accompanies
these data (Supplementary Material \S S8); the corpus is not interpreted here.

\subsection{The ceremonial context}

Participants reported extensive but not recent ceremonial histories. Among the 67 answering, the median
number of lifetime ceremonies attended was 6 (IQR 3--12, range 1--250) against a median of 1 in the past 12
months (IQR 0--2), and 30 of the 66 answering the recency item (45\%) reported their last ceremony more
than 12 months earlier.

Across all ceremonies attended, the substances most often reported by the 67 participants answering were
magic mushrooms (50, 75\%), ayahuasca (43, 64\%), tobacco (38, 57\%), and cacao (23, 34\%), followed by an
``other'' category (18, 27\%) that only 3 respondents specified in its free-text follow-up, leaving it
uninterpretable, and then DMT (17, 25\%), LSD (13, 19\%), MDMA (10, 15\%), and ketamine (4, 6\%)
(Figure~\ref{fig:subs}). For the most recent ceremony, the individually reportable groups were magic
mushrooms ($n=31$) and ayahuasca ($n=23$), every remaining substance being pooled into a combined group
($n=11$) under a prespecified rule.

The most recent ceremony was typically a discrete, multi-hour, outdoor event: a single ceremony rather than
a retreat for 48 of the 65 answering (74\%), a natural setting for 37 of 67 (55\%), and most often either
4--6 hours or an overnight event of roughly 12 hours or more (Table~\ref{tab:context}).

\subsection{Preparation, screening, and prior knowledge}

Preparation was near-universal among those answering (Table~\ref{tab:safety}): of 67, 55 (82\%) prepared
and 6 (9\%) sometimes did, against 6 (9\%) reporting none. This imbalance is why the four-arm design was
abandoned and it constrains every preparation analysis below. Participant-observable facilitator practices
were common but not uniform: 54 of 67 (81\%) knew about the facilitator's background beforehand, but only
39 of the 66 answering the screening item (59\%) reported being asked screening questions, against 15
(23\%) who were not and 7 (11\%) who were unsure. Disclosure outside the ceremonial context was uncommon:
8 of the 65 answering (12\%) had spoken to a health professional about their participation beforehand.

A participant-side index counting these three practices (range 0--3) was constructed for the 54 with
complete data on all three: 35 (65\%) reported all three and 16 (30\%) two, and \emph{none} reported zero.
Internal consistency was low ($\alpha = 0.14$, 95\% CI $-0.36$--0.47), consistent both with the items
capturing distinct protective practices rather than one latent construct and with their differing
referents, since preparation is asked as a general practice while the other two ask about the most recent
ceremony. The score is read as a count of practices present rather than as a validated scale.

\subsection{Experience and perceived safety}

Reported satisfaction was overwhelmingly positive. Of the 72 respondents identifying as participants, 64
answered the satisfaction item---63 giving a numeric score and 1 declining---while 3 were never routed to
the participant section and 5 did not respond. The distribution was strongly ceiling-weighted with a median
of 6 and the mode at the scale maximum (Figure~\ref{fig:good}). Perceived in-ceremony comfort and safety
was similarly high and undifferentiated, returning a median of 3 in every subgroup tested and differing
neither by prior knowledge of the facilitator's background nor by whether screening questions were asked.
Trust in the facilitator was the one participant-side comparison differing at the uncorrected
level---median 4 (IQR 3--4) among the 54 who knew the facilitator's background beforehand against 3 among
both smaller groups ($H=8.15$, $p=0.017$, $\varepsilon^2=0.10$)---and did not survive correction
($q=0.102$). Satisfaction did not differ significantly by most-recent-ceremony substance ($H=4.98$,
$p=0.083$, $q=0.166$). Full test tables are in Supplementary Tables~S4--S7.

The safety domain produced the sample's most consequential descriptive finding. Of the 67 participants
answering the lifetime non-consensual contact item, 61 (91\%) reported no such contact, 5 (7\%) reported
that it occurred, and 1 was unsure; the within-sample proportion was 7.5\% (Wilson 95\% CI 3.2--16.3). The
parallel most-recent-ceremony item, answered by 65, returned 2 affirmative and 62 negative, with 1
declining: of the 5 reporting the experience ever, 2 placed it at the most recent ceremony and 3 did not.
That bears directly on what follows, because the facilitator-practice measures crossed against the lifetime
item all describe the most recent ceremony, so for those 3 respondents exposure and outcome refer to
different events. Non-consensual contact was associated with prior knowledge of the facilitator's
background at the uncorrected level (permutation $\chi^2 = 18.90$, $p=0.047$, $V=0.38$, $n=67$) but not
with screening (permutation $\chi^2 = 13.64$, $p=0.094$, $V=0.32$, $n=66$); neither survived correction
(both $q=0.236$), nor did the comparison by substance ($p=0.486$) or by the harm-reduction index
($p=0.641$). The research team elected to publish these cross-tabulations rather than withhold them:
withholding would remove the only evidence this study can offer on whether such reports cluster with any
measured practice, and aggregate reporting of an anonymous survey that collected no identifiers leaves no
respondent identifiable. Every result here is hypothesis-generating, rests on the referent mismatch just
described, and is powered neither to confirm nor to rule out an association. One measurement observation
belongs alongside them: an item on perceived positive life impact had zero variance, all 67 respondents
answering selecting the same option.

\subsection{Facilitator practice}

Eight respondents identified as facilitators, and 7 of the 8 also identified as participants. Every
facilitator who answered a substantive facilitator item is one of those 7, so this section and the
participant results above describe overlapping people rather than independent samples. Cells are small and
no comparison at this base is meaningful. Of the 7 answering, 6 gave personal experience and self-study as
their preparation for the role and \emph{none} reported formal certification or academic coursework; on
emergency planning, 4 described an informal or unwritten approach and 3 said it depends on the situation,
with none selecting a formal written plan. All 7 answered the crisis-signs item: disorientation, feeling
overwhelmed, fainting, difficulty breathing, and panic attacks were each reported by 3, and signs of
self-harm or suicidal thoughts by 2. Full distributions, including screening and distress-response
practices, and the undefined odds ratios this branch returns, are in Supplementary Tables~S2--S3.

\subsection{Preparation and reported satisfaction}

Among the 63 participants giving a numeric satisfaction score, 58 reported some preparation (median 6, IQR
5--6) and 5 reported none (median 4, IQR 3--5; Figure~\ref{fig:prep}); a sixth respondent reporting no
preparation did not answer the outcome, which is why this base is 5 rather than the 6 in
Table~\ref{tab:safety}. The Mann--Whitney comparison specified in the analysis plan returned $U=218.0$,
$p=0.041$, $r_{rb}=0.50$, and a common-language effect size of 0.75; a tie-aware permutation test gave
$p=0.024$. The structural weakness of this comparison precedes its result: the no-preparation group
contains five respondents. The prespecified three-level sensitivity analysis was non-significant ($H=4.24$,
$p=0.120$, $q=0.359$), the practice index was unassociated with satisfaction ($H=2.09$, $p=0.352$,
$q=0.422$), and an exploratory proportional-odds model adjusting for age and facilitator trust placed
preparation short of conventional significance (OR 5.21, 95\% CI 0.84--32.41, $p=0.077$) while trust
reached it (OR 2.20, 95\% CI 1.06--4.53, $p=0.033$), at roughly 12 cases per estimated parameter
(Supplementary Table~S8). No test in any corrected family was significant at $q<0.05$.

\section{Discussion}

\subsection{A structurally informal practice landscape}

The most stable picture this sample supports is a practice organized around informal, relationally
transmitted expertise rather than credentialed or documented protocol. Among the facilitators who
self-reported, training came from personal experience, self-study, mentorship, and family tradition; none
reported formal certification or academic coursework, and emergency planning was uniformly informal or
conditional. This sits at some distance from the structured preparation, screening, and difficult-experience
protocols that harm-reduction frameworks for psychedelic settings assume
\citep{gorman2021psychedelic,pilecki2021ethical,dutton2025harm}, is consistent with survey evidence that
people who use psychedelics support clearer practice standards \citep{kruger2023preferences}, and is a
reminder that in community settings the containing structure is the setting itself rather than a protocol
\citep{hartogsohn2017constructing}.

On the participant side the informality reads differently. Preparation was near-universal among those
answering, but roughly a quarter reported that no screening questions were asked and a further group was
uncertain. The asymmetry is worth separating: preparation is something participants largely do, whether
self-directed or instructed, whereas screening must be done \emph{to} them, and was substantially less
consistent here. If that asymmetry holds in a larger sample it would locate the intervention target on the
facilitator side---but this sample cannot establish whether it holds, and eight facilitators cannot
characterize screening practice.

The demographic profile---middle-aged, university-educated, economically comfortable---resembles that of
ceremony-attender cohorts described elsewhere, though at the older and more educated end of them. Its mean
age of 42.4 years and 78\% holding a bachelor's degree or above compare with 37.6 years and 49\%
degree-holding among United States neo-shamanic ayahuasca ceremony attenders \citep{pagni2025longterm},
37.0 years among ceremonial ayahuasca retreat attenders in the Peruvian Amazon
\citep{ruffell2021ceremonial}, and 40 years with 59\% degree-holding in a large international naturalistic
psilocybin cohort \citep{nayak2023naturalistic}; the 49\% here identifying as women falls inside the
44--61\% those three report, as it does for international psychedelic-survey samples
\citep{kopra2023investigation,robinson2026field} and this team's psilocybin-specific Puerto Rico survey
\citep{velezrodriguez2026psilocybin}. That convergence is best read as evidence about who responds to online
psychedelic surveys, not about who attends ceremonies \citep{bethlehem2010selection}.

\subsection{A ceiling that limits what satisfaction can reveal}

Reported satisfaction clustered at the top of its scale, and a perceived-life-impact item had no variance
at all. This is a substantive finding---these respondents overwhelmingly described their most recent
ceremony positively, consistent with naturalistic cohorts reporting predominantly positive appraisals
\citep{nayak2023naturalistic,ruffell2021ceremonial}---but it also constrains the analysis severely. An
outcome with almost no variance in the lower range cannot discriminate between groups, and the null results
across the corrected families are consistent with that constraint rather than with the absence of any
association. Ceiling behavior is a known property of self-report scales for hallucinogen experience rather
than a peculiarity of this instrument \citep{bouso2016measuring}. Satisfaction, measured this way, is close
to exhausted as an outcome here; future work needs measures with headroom in both directions before group
comparisons on experience quality can be informative.

\subsection{Safety signals the design can identify but not quantify}

Five of 67 participants answering reported non-consensual physical contact at some point, two of them at
their most recent ceremony, and facilitators reported having witnessed crisis signs at a base of seven at
which no rate is estimable. These come from different reporters---none of the five reporting
non-consensual contact is a facilitator---and cannot be combined. That such reports appear at all is
consistent with survey-based evidence on boundary violations in naturalistic psychedelic settings
\citep{kruger2025psychedelic,mcnamee2023studying} and with the community harm-reduction literature that
developed in response \citep{peluso2020reflections}, and the crisis signs described are of a kind
documented at non-trivial rates in large naturalistic samples
\citep{carbonaro2016survey,bouso2022adverse}. The item used here was deliberately broad, covering
unexpected touch and uncertain or unremembered consent, so the figure should not be read against a narrower
legal or clinical standard.

The study establishes that these events occur among respondents and that an instrument of this kind can
capture them; it cannot establish how often they occur, in whom, or under what conditions. The
cross-tabulations against participant-side proxies pair a lifetime outcome with single-ceremony exposures,
are powered neither to confirm nor to rule out an association, and did not survive correction, so treating
any of them as an actionable signal would overstate them considerably. The safety domain justified its
inclusion by identifying a measurement priority, not by producing an estimate.

\subsection{The preparation finding as a design lesson}

The preparation-to-satisfaction comparison is this study's pre-specified comparison and its most fragile
result, and those two facts should be held together rather than traded off. The direction was positive and
the effect sizes were not small, but the no-preparation group contained five respondents answering the
outcome, the three-level sensitivity analysis was non-significant, and the adjusted model placed
preparation short of significance with a confidence interval spanning nearly two orders of magnitude.
Nothing here establishes the finding.

The durable contribution is diagnostic: the preparation split was unbalanced enough to dissolve the
original four-arm design and remains the binding constraint on the primary analysis, so any successor study
will need either a sampling strategy that reaches unprepared participants in usable numbers or a question
that does not depend on that contrast. Reporting the deviation rather than absorbing it quietly is what
makes the constraint visible to the next design.

\subsection{Limitations}

The sample was non-probability, unweighted, and self-selected into an online survey; the concentration in
higher education categories and favorable economic self-ratings describe who responded, not who
participates, and nothing here supports a statement about how common ceremonial entheogen use is in Puerto
Rico. Roughly 6\% of those answering the residency item lived outside Puerto Rico, further tempering any
Puerto Rico--specific reading, and the eligibility criteria were self-assessed rather than enforced.

Role subsamples were small and unevenly sized, and the facilitator branch of eight supports no rate and no
association test. Because 7 of those 8 also answered as participants, the facilitator-side and
participant-side statistics are not independent samples and should not be read as two separate groups. The
reports are nonetheless unlinked at the ceremony level, which is not a minor inconvenience: it eliminated
the protocol's integration-proxy comparison, its perceived-safety-by-training analysis, and its mixed
harm-reduction composite, and no statement in this paper connects what facilitators reported doing to what
participants reported experiencing.

Two referent mismatches limit the associations reported here: the safety cross-tabulations pair a lifetime
outcome with most-recent-ceremony exposures, and the primary comparison pairs a most-recent-ceremony
outcome with a general-practice exposure. All measures were single-timepoint self-report, and for most
respondents the most recent ceremony was distant---30 of 66 answering reported more than 12 months---so
recall error is plausible and unmeasurable. The instrument contains no integration measure, is unvalidated,
and produced mixed first reliability evidence (Supplementary Table~S9) alongside two documented skip-logic
errors (Supplementary Material \S S6).

\subsection{Conclusion}

Within a self-selected sample of 73 adults, ceremonial entheogen use in Puerto Rico was described by a
middle-aged, university-educated, overwhelmingly Puerto Rico--resident group whose most recent ceremonies
were typically single, multi-hour, outdoor events involving psilocybin mushrooms or ayahuasca, described in
strongly positive terms. Preparation was near-universal while pre-ceremony screening was not, and the small
facilitator subsample described training and emergency planning as informal rather than credentialed. Five
of 67 participants answering reported non-consensual physical contact, and facilitators reported witnessing
crisis events at a base too small for any rate to be estimated. These findings describe this sample alone,
support no estimate of prevalence and no inference about what any practice produces, and take their
principal value from identifying which measures, sampling strategies, and instrument corrections a properly
powered successor study would require.

\section*{Author contributions}
% UNCONFIRMED — drafted from the roles described in the study protocol. Requires author sign-off.
YOOO conceived the study, led protocol development and ethics submission, and supervised the project.
YOOO, JCVR, JMHT, JMT, PRF, CR, ANGR, and AIRM contributed to instrument development and review. JCVR led
data curation, formal analysis, and manuscript drafting. JMHT contributed to methodology and analysis
review. JMT, PRF, and CR contributed to interpretation and clinical/community context. ANGR and AIRM
contributed to recruitment and data collection. All authors reviewed and approved the final manuscript.

\section*{Funding}
% Confirmed by the corresponding author, 2026-08-13: the study was unfunded.
This research received no specific grant from any funding agency in the public, commercial, or
not-for-profit sectors.

\section*{Competing interests}
% Author decision, 2026-08-13: disclose the organizational affiliations rather than declare none.
% Still requires a per-author disclosure before submission — UNCONFIRMED at the individual level.
Several authors are affiliated with organizations that conduct, advocate for, or provide services related
to psychedelic and entheogen research and practice, including the Puerto Rico Institute for Psychedelic
Science, Medicine \& Awareness and the Colectivo Psicod\'elico de Puerto Rico. These affiliations are
disclosed as potential non-financial competing interests. No author received financial support from, or
holds a financial interest in, any commercial entity with an interest in the subject matter of this
article.

\section*{Data availability}
% Author decision, 2026-08-13: commit to a public repository deposit rather than release on request.
The de-identified survey data and the complete analysis code that generates every result reported here will
be deposited in a public repository on publication, and the digital object identifier will be recorded
here. The data contain no direct identifiers by design. Until deposit, both are available from the
corresponding author.

\section*{Generative AI use declaration}
% UNCONFIRMED — verify scope before submission.
Large language model tooling (Anthropic Claude) was used to assist with statistical analysis code
development and with drafting and editing of this manuscript. All analyses were specified by the authors,
all outputs were verified against the source data by the authors, and the authors take full responsibility
for the content of this article.

\bibliography{../references/references}

\clearpage

\section*{Tables and Figures}

\begin{table}[htbp]
\centering
\caption{Sociodemographic characteristics of the analytic sample ($N=73$). Bases differ by item and are
stated with each; counts are exact and percentages are of respondents answering.}
\label{tab:demo}
\setstretch{1.1}
\begin{tabular}{p{3.6cm}p{6.4cm}rr}
\toprule
Characteristic & Category & $n$ & \% answering \\
\midrule
Age, years ($n=72$) & Mean 42.4 (SD 11.9); median 41 (IQR 34--50); range 21--71 & & \\
\addlinespace
Gender ($n=71$)     & Woman & 35 & 49 \\
                    & Man & 30 & 42 \\
                    & Non-binary & 4 & 6 \\
                    & Prefer not to answer & 2 & 3 \\
\addlinespace
Education ($n=72$)  & Bachelor's degree & 34 & 47 \\
                    & Master's degree & 19 & 26 \\
                    & Associate's degree & 7 & 10 \\
                    & High school diploma / GED & 6 & 8 \\
                    & Doctorate & 3 & 4 \\
                    & Prefer not to answer & 3 & 4 \\
\addlinespace
Economic situation ($n=72$) & Comfortable & 44 & 61 \\
                    & Tight & 20 & 28 \\
                    & Difficult & 6 & 8 \\
                    & Prefer not to answer & 2 & 3 \\
\addlinespace
Country of residence ($n=72$) & Puerto Rico & 68 & 94 \\
                    & United States & 3 & 4 \\
                    & Spain & 1 & 1 \\
\bottomrule
\end{tabular}
\end{table}

\begin{table}[htbp]
\centering
\caption{Most-recent-ceremony context, participant-reported. Participant branch only: 72 respondents
identified as participants, of whom 3 were not routed to these items by the continuation gate described in
Methods and appear as structural skips. Item nonresponse is the residual between the per-item base and 69.
Counts are exact; percentages are of respondents answering the item.}
\label{tab:context}
\setstretch{1.0}
\begin{tabular}{p{4.2cm}p{5.8cm}rr}
\toprule
Item & Category & $n$ & \% answering \\
\midrule
Time since last ceremony & More than 12 months ago & 30 & 45 \\
($n=66$ answering)                & 2 to 6 months ago & 20 & 30 \\
                                  & 7 to 12 months ago & 11 & 17 \\
                                  & During the past month & 4 & 6 \\
                                  & Prefer not to answer & 1 & 2 \\
                                  & \emph{Structural skip (not routed)} & 3 & --- \\
                                  & \emph{Item nonresponse} & 3 & --- \\
\addlinespace
Place ($n=67$ answering)   & Natural setting (forest, beach, mountain) & 37 & 55 \\
                 & Private home & 15 & 22 \\
                 & Retreat center or wellness facility & 9 & 13 \\
                 & Other & 3 & 4 \\
                 & Urban or residential space used informally & 2 & 3 \\
                 & Prefer not to answer & 1 & 1 \\
                 & \emph{Structural skip (not routed)} & 3 & --- \\
                 & \emph{Item nonresponse} & 2 & --- \\
\addlinespace
Format ($n=65$ answering)  & Single ceremony & 48 & 74 \\
                 & Retreat or series of ceremonies & 13 & 20 \\
                 & Prefer not to answer & 3 & 5 \\
                 & Not sure & 1 & 2 \\
                 & \emph{Structural skip (not routed)} & 3 & --- \\
                 & \emph{Item nonresponse} & 4 & --- \\
\addlinespace
Duration ($n=65$ answering) & 4 to 6 hours & 19 & 29 \\
                 & Overnight ($\sim$12+ hours) & 19 & 29 \\
                 & 7 to 12 hours & 14 & 22 \\
                 & 1 to 3 hours & 7 & 11 \\
                 & 2 to 3 days & 4 & 6 \\
                 & 4 to 6 days & 1 & 2 \\
                 & Less than 1 hour & 1 & 2 \\
                 & \emph{Structural skip (not routed)} & 3 & --- \\
                 & \emph{Item nonresponse} & 4 & --- \\
\bottomrule
\end{tabular}
\end{table}

\setstretch{1.1}
\begin{longtable}{p{5.2cm}p{4.8cm}rr}
\caption{Participant-reported protective practices and safety outcomes. Bases are per item and differ
because each is asked under its own skip condition; the 3 structural skips are the dual-role respondents
not routed to the participant section. Counts are exact throughout. The two non-consensual contact items
differ in referent: the first asks about any ceremony ever, the second about the most recent ceremony only.}
\label{tab:safety} \\
\toprule
Item & Category & $n$ & \% answering \\
\midrule
\endfirsthead
\toprule
Item & Category & $n$ & \% answering \\
\midrule
\endhead
Self-directed preparation ($n=67$) & Yes & 55 & 82 \\
                & Sometimes & 6 & 9 \\
                & No & 6 & 9 \\
                & \emph{Structural skip / item nonresponse} & 3 / 2 & --- \\
\addlinespace
Knew facilitator's background ($n=67$) & Yes & 54 & 81 \\
                & A little & 8 & 12 \\
                & No & 4 & 6 \\
                & Prefer not to answer & 1 & 1 \\
                & \emph{Structural skip / item nonresponse} & 3 / 2 & --- \\
\addlinespace
Facilitator asked screening questions ($n=66$) & Yes & 39 & 59 \\
                & No & 15 & 23 \\
                & Not sure & 7 & 11 \\
                & Prefer not to answer & 5 & 8 \\
                & \emph{Structural skip / item nonresponse} & 3 / 3 & --- \\
\addlinespace
Spoke to a health professional beforehand ($n=65$) & No & 54 & 83 \\
                & Yes & 8 & 12 \\
                & Prefer not to answer & 2 & 3 \\
                & Not sure & 1 & 2 \\
                & \emph{Structural skip / item nonresponse} & 3 / 4 & --- \\
\addlinespace
Non-consensual physical contact, ever ($n=67$) & No & 61 & 91 \\
                & Yes & 5 & 7 \\
                & Not sure & 1 & 1 \\
                & \emph{Structural skip / item nonresponse} & 3 / 2 & --- \\
\addlinespace
Non-consensual physical contact, most recent ceremony ($n=65$) & No & 62 & 95 \\
                & Yes & 2 & 3 \\
                & Prefer not to answer & 1 & 2 \\
                & \emph{Structural skip / item nonresponse} & 3 / 4 & --- \\
\addlinespace
Harm-reduction practice index ($n=54$ complete cases) & All three practices & 35 & 65 \\
                & Two practices & 16 & 30 \\
                & One practice & 3 & 6 \\
                & No practices & 0 & 0 \\
\bottomrule
\end{longtable}
\setstretch{1.0}

\begin{figure}[htbp]
\centering
\includegraphics[width=0.94\textwidth]{../figures/fig1_flow.png}
\caption{Sample construction and role structure. Of 100 submissions, 13 were abandoned before the role
fork and 14 of the remaining 87 screened out by answering ``no'' to both role items, leaving an analytic
sample of $N=73$. Roles are not mutually exclusive: the 7 dual-role respondents enter both branches, giving
72 participant-side and 8 facilitator-side reports, and 3 of those 7 were stopped by the participant
continuation gate, so 69 respondents were actually routed to the participant items. The branches are
unlinked at the ceremony level --- no identifier connects a participant to the facilitator who ran their
ceremony, and the ceremonies a dual-role respondent led are different events from those they attended ---
so no participant--facilitator comparison is possible.}
\label{fig:flow}
\end{figure}

\begin{figure}[htbp]
\centering
\includegraphics[width=\textwidth]{../figures/fig2_substances.png}
\caption{Substances used across all ceremonies attended, participant-reported (select-multiple; base $n=67$
participants answering). Bars are labeled with the count and the percentage of that base; because
respondents could select more than one substance, percentages need not sum to 100\%. The ``other'' bar is
drawn in grey because its free-text follow-up was completed by only 3 respondents, leaving the category
uninterpretable. Distinct from the facilitator-led substance variable and from the most-recent-ceremony
substance variable.}
\label{fig:subs}
\end{figure}

\begin{figure}[htbp]
\centering
\includegraphics[width=\textwidth]{../figures/fig3_ceremony_good.png}
\caption{Distribution of reported satisfaction with the most recent ceremony (0 = very unpleasant to 6 =
very pleasant). Of the 72 respondents identifying as participants, 63 gave a numeric score and form the
base; the remaining 9 comprise 1 ``prefer not to answer'', 3 structural skips at the continuation gate, and
5 item nonresponses. Median 6, with 51 of 63 (81\%) answering 5 or 6 --- the ceiling that limits how much
variance any predictor can explain.}
\label{fig:good}
\end{figure}

\begin{figure}[htbp]
\centering
\includegraphics[width=\textwidth]{../figures/fig4_prep_outcome.png}
\caption{Reported satisfaction by self-directed preparation status. Any preparation, $n=58$ answering the
outcome (median 6, IQR 5--6); no preparation, $n=5$ (median 4, IQR 3--5). Marker area is proportional to the
number of respondents giving each score, since scores tie heavily at the ceiling. The no-preparation group
is drawn as its five raw points and a median tick rather than a box, which at that $n$ would imply
distributional precision the data cannot support. \textbf{Exploratory:} this comparison is suggestive, not
confirmatory.}
\label{fig:prep}
\end{figure}

\end{document}
